\section{Terminology and Witness Register}
\label{app:terminology-witnesses}

\subsection{Theological Terminology and Translation Collisions}

This glossary records the preferred English terms used across this treatise, their Hebrew, Greek, and Latin originals, and the translation collisions that arise between Greek patristic texts, Latin scholasticism, and nineteenth-century English versions (above all John Parker's 1899 rendering of the \emph{Celestial Hierarchy}).

{\footnotesize
\setlength{\tabcolsep}{2.5pt}
\renewcommand{\arraystretch}{1.15}
\begin{longtable}{>{\bfseries\raggedright\arraybackslash}p{0.17\linewidth}>{\raggedright\arraybackslash}p{0.18\linewidth}>{\raggedright\arraybackslash}p{0.31\linewidth}>{\raggedright\arraybackslash}p{0.28\linewidth}}
\toprule
\textbf{Preferred Term} & \textbf{Source Terms} & \textbf{Meaning} & \textbf{Collisions and Usage} \\
\midrule
\endfirsthead
\toprule
\textbf{Preferred Term} & \textbf{Source Terms} & \textbf{Meaning} & \textbf{Collisions and Usage} \\
\midrule
\endhead
\bottomrule
\endfoot
Angel & \gk{mal'ak} / \gk{angelos} / \emph{angelus} & Name of office (``messenger''), not of nature (\emph{spiritus}); also proper name of the ninth order (\CH{5, 9}; Gregory, \emph{Hom.}~34.8). & Distinguished from the lowest choir specifically and from Christ as ``Angel of Great Counsel'' (Is~9:5~LXX). \\
\addlinespace[2pt]
Demon / Devil & \gk{satan} / \gk{diabolos}, \gk{daimonion} / \emph{daemon} & Exclusively a fallen, wicked spiritual creature. & Breaks with Greek neutral/benevolent \gk{daim\={o}n} (Hesiod, Plato, Apuleius); follows Ps~95:5~LXX and Augustine (\emph{Civ.\ Dei}~IX.19). \\
\addlinespace[2pt]
Dominations & \gk{Kyriot\={e}tes} / \emph{Dominationes} & Fourth order (Triad~II.1): free lordship directing ministries. & Parker~(1899) renders ``Lordships''; older versions also use ``Dominions.'' \\
\addlinespace[2pt]
Virtues & \gk{Dynameis} / \emph{Virtutes} & Fifth order in Dionysius (Triad~II.2); seventh in Gregory (\emph{Hom.}~34.7): virile strength executing great works. & \textbf{Collision:} Parker~(1899) translates \gk{Dynameis} as ``Powers,'' colliding with Latin \emph{Potestates}. \\
\addlinespace[2pt]
Powers & \gk{Exousiai} / \emph{Potestates} & Sixth order (Triad~II.3 in Dionysius and Gregory \emph{Hom.}~34.7): ordering and restraining hostile forces. & \textbf{Collision:} Parker~(1899) translates \gk{Exousiai} as ``Authorities.'' \\
\addlinespace[2pt]
Principalities & \gk{Archai} / \emph{Principatus} & Seventh order in Dionysius (Triad~III.1); fifth in Gregory (\emph{Hom.}~34.7): leadership over nations and public ministry. & Rendered ``Princedoms'' in some 19th-century patristic translations. \\
\addlinespace[2pt]
Subsisting form & \emph{forma subsistens}, \emph{substantia separata} & Pure immaterial form without matter, composed of essence (potency) and received \emph{esse} (act) (\ST{I, q.~50, a.~2}). & Contrasted with Bonaventure's universal hylomorphism (``spiritual matter''). \\
\addlinespace[2pt]
Aevum & \emph{aevum} & Measure of angelic substantial duration, having a beginning in creation without substantial change (\ST{I, q.~10, a.~5}). & Distinct from discrete angelic time measuring successive acts (\ST{I, q.~53, a.~3}). \\
\addlinespace[2pt]
Virtual contact & \emph{contactus virtutis} & Presence of an angel in a place by applying active spiritual power rather than dimensional extension (\ST{I, q.~52, a.~1}). & Contrasted with bodily dimensive contact and Scotus's definitive presence. \\
\addlinespace[2pt]
Innate intelligible species & \emph{species intelligibiles concreatae} & Intellectual similitudes impressed by God at creation, enabling knowledge without sensory abstraction (\ST{I, q.~55, a.~2}). & Distinct from human abstracted species and from the beatific vision of God's essence. \\
\addlinespace[2pt]
Morning / evening knowledge & \emph{cognitio matutina / vespertina} & Knowing creatures in the Word (morning) vs.\ in their proper created nature (evening) (Augustine; \ST{I, q.~58, aa.~6--7}). & Contrasted with demonic ``night'' knowledge refusing reference to God. \\
\addlinespace[2pt]
Illumination vs.\ Speech & \emph{illuminatio} vs.\ \emph{locutio} & Illumination conveys divine truth downward (\ST{I, q.~106}); speech voluntarily directs an interior concept bidirectionally and privately (\ST{I, q.~107}). & Lower angels speak to higher angels without illuminating them. \\
\addlinespace[2pt]
Latria vs.\ Dulia & \gk{latreia} vs.\ \gk{douleia} & \emph{Latria} is adoration due to God alone; \emph{dulia} is lawful honor and invocation offered to angels and saints (\emph{hyperdulia} to Mary). & Guard against pagan theurgy (Augustine, \emph{Civ.\ Dei}~X) and Laodicea~c.~35 angelolatry. \\
\addlinespace[2pt]
Apocryphal & \gk{apokrypha} & Extra-canonical Jewish writings (\emph{1~Enoch}, \emph{Jubilees}, \emph{Testaments}) outside the biblical canon (Augustine, \emph{Civ.\ Dei}~XV.23). & Distinct from deuterocanonical Scripture (e.g., Tobit), which is canonical. \\
\addlinespace[2pt]
Pious tradition & \emph{pia traditio} & Devotional practices and saintly accounts received in Catholic spirituality without magisterial definition. & Presented affirmatively as received, distinct from defined dogma. \\
\addlinespace[2pt]
\emph{Potentia absoluta / ordinata} & \emph{potentia absoluta / ordinata} & What God can do by simple omnipotence vs.\ what conforms to His ordained wisdom (\emph{De malo}, q.~16, a.~5~ad~13). & Clarifies that demonic obstinacy does not limit divine omnipotence. \\
\end{longtable}
}

\subsection{Primary Witness Register}

The summary register below categorizes the principal editions and textual controls consulted in preparing this treatise, distinguishing complete public-domain editions inspected directly from modernized revisions, verified passage records, and works summarized through secondary controls or official summaries without full electronic editions.

{\footnotesize
\setlength{\tabcolsep}{3pt}
\renewcommand{\arraystretch}{1.15}
\begin{longtable}{>{\bfseries\raggedright\arraybackslash}p{0.19\linewidth}>{\raggedright\arraybackslash}p{0.35\linewidth}>{\raggedright\arraybackslash}p{0.39\linewidth}}
\toprule
\textbf{Corpus Family} & \textbf{Primary Editions Inspected} & \textbf{Textual and Verification Ceiling} \\
\midrule
\endfirsthead
\toprule
\textbf{Corpus Family} & \textbf{Primary Editions Inspected} & \textbf{Textual and Verification Ceiling} \\
\midrule
\endhead
\bottomrule
\endfoot
Sacred Scripture & Douay-Rheims Challoner Bible; Clementine Vulgate; Brenton English Septuagint; Swete Greek Psalms; versification concordances & Public-domain verbatim verse control; Swete covers the Psalter in Greek (Deuteronomy, Isaias, Daniel, and Job read in Brenton's Septuagint); critical academic Septuagint editions remain restricted. \\
\addlinespace[2pt]
St.~Thomas Aquinas & \emph{ST}~I qq.~50--64 (English Dominican 1920, Gutenberg revision, and 1889 Leonine Latin); \emph{ST}~I qq.~106--114 and III qq.~8, 30 (English Dominican); \emph{De malo} q.~16, \emph{De veritate} qq.~8--9, \emph{De ente} chs.~4--5 (Corpus Thomisticum Latin); \emph{Super Eph./Col.} (Dessain 1857 OCR) & Three electronic printing errors corrected against the Leonine Latin; Dessain 1857 OCR treated strictly as a finding aid; \emph{SCG}, \emph{De spir.\ creat.}, \emph{De subst.\ sep.}, and \emph{In II Sent.} not consulted in full electronic editions. \\
\addlinespace[2pt]
Pseudo-Dionysius & Parker~(1899) complete English \emph{CH}; Migne \emph{PG}~3 complete Greek-Latin text (cols.~119--340; plus \emph{DN}~IV.23, \emph{EH}) & All fifteen chapters and \emph{PG}~3 column spans verified directly in full text rather than page-image facsimiles. \\
\addlinespace[2pt]
Pre-Christian Writings & Brenton Septuagint; Swete Greek Psalms; verified passage records for Plato \emph{Symposium} 202d--203a and \emph{Jubilees} 10:1--11; direct quotations inside Augustine \emph{Civ.\ Dei} VIII--X, XV.23 and Justin \emph{2~Apol.}~5 & \emph{1~Enoch}, \emph{Testaments}, Philo, Hesiod, Plutarch, Apuleius, Plotinus, Porphyry, and Iamblichus were not consulted in full standalone editions and are reported strictly through patristic controls and verified passage records. \\
\addlinespace[2pt]
Church Fathers & Unrevised ANF~1 and NPNF~1--2 printings (Clement of Rome, Justin, Irenaeus, Athanasius, Cyril of Jerusalem, Basil, Nazianzen, Nyssa, Chrysostom, Hilary, Ambrose, Jerome, Leo, Damascene); Dods~1871 \emph{City of God}; Library of the Fathers \emph{Moralia}; Latin Gregory \emph{Hom.\ in Ev.}; Migne \emph{PG}/\emph{PL} controls & Two texts derive from modernized New Advent revisions (Origen \emph{In Matt.}~XIII.26--29; Augustine \emph{De Trin.}~III.4--9) and are paraphrased rather than presented as unrevised nineteenth-century printings. \\
\addlinespace[2pt]
Magisterium & Hefele (Laodicea~c.~35; Constantinople~543/553; Roman Synod~745 summary); Denzinger Latin (Lateran~IV DS~800--801); 1860 \emph{Bullarium} (Florence); Trent \emph{Canones} and 1566 \emph{Roman Catechism}; Vatican~I \emph{Dei Filius} (Latin and English) & Council of Vienne (\emph{Ad nostrum qui}, 1312) and modern Vatican acts (\emph{Humani generis}, \emph{Ad caeli Reginam}, \emph{Lumen gentium}, CDF~1975, CDF~1992, CCC, CDWDS~2001) are summarized accurately without reproducing restricted or uninspected verbatim text. \\
\addlinespace[2pt]
Later Schools and Byzantine & Bonaventure \emph{Breviloquium} II.6--8 (Latin) and \emph{In II Sent.} d.~3 (1885 Quaracchi vol.~2, supplying Peter Lombard's Book~II text and Quaracchi \emph{Scholia} on Scotus and Cajetan); John of Damascus \emph{DFO} II.2--4 & Direct standalone editions of Scotus (\emph{Ordinatio}), Su\'arez (\emph{De angelis}), Maximus (\emph{Mystagogia}), Theodore the Studite, and Gregory Palamas were not consulted in full and are summarized without fabricated quotations. \\
\addlinespace[2pt]
Liturgy, Prayers, and Saints & 1889 Descl\'ee and 1945 Lasance Missals; 1962 Roman Missal and 1981 \emph{Ordo Lectionum Missae} calendar records; 1910 \emph{Raccolta} and \emph{ASS}~26 (received prayers); Bernard (\emph{Qui habitat}, \emph{Missus est}), Teresa (\emph{Life}), Francis de Sales (\emph{Introduction}), Alphonsus (\emph{Preparation}, \emph{Incarnation}), Th\'er\`ese (\emph{Story of a Soul}) & Postconciliar ICEL English translations and postconciliar Offerings/Postcommunions are withheld for rights; the Byzantine 8~November Menaion supplies calendar facts only (no prayer text). \\
\end{longtable}
}
